\documentclass[10pt,a4paper]{amsart}
\usepackage[margin=19mm]{geometry}
\usepackage{amsmath,amssymb,mathtools}
\usepackage[scr=rsfs,cal=euler]{mathalfa}
\usepackage{xfrac}
\usepackage[unicode,hidelinks,bookmarks=false]{hyperref}
\newcommand{\ie}{i.e.\ }
\newcommand{\cf}{cf.\ }
\newcommand{\resp}{respectively\ }
\newcommand{\Subsection}{\subsection}
\providecommand{\nobreakdash}{\nobreak\hbox{-}}
\setlength{\emergencystretch}{2em}
\multlinegap0pt
\usepackage{amssymb}
\usepackage{tikz}
\usepackage{braket}
\newtheorem{prop}{Proposition}
\title{Equivariant Quantum Cohomology of Grassmannians via the Clifford algebra}
\author{Christian Korff, Mikhail Vasilev}
\date{Source: arXiv:2606.06352v1 (CC BY 4.0)}
\begin{document}
\maketitle

\noindent\emph{Source and licence.} Equivariant Quantum Cohomology of Grassmannians via the Clifford algebra. Version arXiv:2606.06352v1 (CC BY 4.0), distributed by arXiv under the Creative Commons Attribution 4.0 International licence, http://creativecommons.org/licenses/by/4.0/, as recorded in the arXiv licence metadata for this exact version and shown on the abstract page https://arxiv.org/abs/2606.06352v1. Full licence notice: \texttt{licenses/arxiv-2606.06352-CC-BY-4.0.txt}.\par\smallskip

\noindent\textit{Editorial scope.} Counted unit: the article's prop ReverseBranching (source0.tex lines 1630--1636). The article's complete original proof of it is delivered with it (source0.tex lines 1637--1665, one \texttt{proof} environment of 29 lines). No mathematics is written, repaired or re-derived: the statement and the proof are the article's own lines, in the article's own notation, and no retained line is substituted or re-flowed.  Retained uncounted as the source-local context the proof needs: lines 1610--1613; lines 1626--1628.  Nothing was omitted from any retained span, so the package carries no omission record.  No retained line contains a citation, so the wrapper carries no bibliography and no bibliography entry.  No retained line contains a figure environment, an image-inclusion command or drawing code, so no image is delivered and the package declares no asset dependency.  Editorial formatting only, in addition: an AMS \texttt{amsart} preamble that restates the article's own declarations and macros used by the retained lines, namely the environments \texttt{prop} on separate counters, together with the standard packages those lines need.  The retained environments print in this document as Proposition 1 in order of appearance, whereas the article prints the same items with the numbering of its own full document.  The complete native source of the article as distributed in the arXiv e-print for this exact version is bundled unchanged as \texttt{sources/202181/source0.tex}.
\begin{equation}
    \label{tablformula}
    s_{\lambda}(x|y) = \sum\limits_{T(\lambda)}   \prod\limits_{\square \in \lambda} (x_{T(\square)} - y_{T(\square) + c(\square)}),
\end{equation}

\begin{equation}\label{BranchingRule}
    s_{\lambda}(x_1, \ldots, x_{k-1},x_k|y) = \sum\limits_{\mu} \prod\limits_{\square \in \lambda / \mu}(x_k - y_{k + c(\square)}) s_{\mu}(x_1, \ldots, x_{k-1}|y),
\end{equation}

\begin{prop}
The factorial Schur polynomials satisfy the reverse branching rule
\begin{equation}\label{ReverseBranching}
    s_{\lambda}(x_1, \ldots, x_{k - 1}|y) = \sum\limits_{\mu}(-1)^{|\lambda/\mu|}\prod\limits_{\square \in \lambda / \mu}(x_k - y_{k + c(\square)}) s_{\mu}(x_1, \ldots, x_{k - 1}, x_k|y),
\end{equation}
where the sum runs over all partitions $\mu \subseteq \lambda$ such that $\lambda / \mu$ is a vertical strip. Note that since factorial Schur polynomials are symmetric with respect to $x$ variables, the branching rule \eqref{BranchingRule} with respect to another $x$ variable is the same.
\end{prop}

\begin{proof}
    We first prove the formula \eqref{ReverseBranching} for complete factorial symmetric polynomials. Indeed, to compute $h_i(x_1, \ldots, x_{k - 1}|y)$ we sum over the semistandard Young tableaux (SSYT) of shape $\lambda = (i)$ filled in with numbers at most $k-1$ via formula \eqref{tablformula}. On the other hand, the expression
    \[
    h_i(x_1, \ldots, x_k|y) - (x_k - y_{k + i - 1})h_{i - 1}(x_1, \ldots, x_k|y)
    \]
    counts all SSYT of the same shape $\lambda = (i)$ filled in with numbers at most $k$ minus all SSYT with at least one $k$ appearing, which is the same as counting all SSYT with at most $k-1$, thus we obtain
    \begin{equation}
    \label{reversecomplete}
    h_i(x_1, \ldots, x_{k - 1}|y) = h_i(x_1, \ldots, x_k|y) - (x_k - y_{k + i - 1})h_{i - 1}(x_1, \ldots, x_k|y).
    \end{equation}
    Now we use the Jacobi-Trudi formula to express the factorial Schur polynomials in terms of factorial complete symmetric polynomials
    \[
    s_{\lambda}(x_1, \ldots, x_{k - 1}|y) = \det(h_{\lambda_{i} - i + j}(x_1, \ldots, x_{k - 1}|\tau^{1 - j}y)).
    \]
    Apply the reverse branching rule \eqref{reversecomplete} for each of the complete factorial symmetric polynomials 
    \begin{multline*}
    h_{\lambda_i - i + j}(x_1, \ldots, x_{k - 1}|\tau^{1-j}y) = h_{\lambda_i - i + j}(x_1, \ldots, x_{k}|\tau^{1-j}y)\\
    - (x_k - y_{k +\lambda_i - i})h_{\lambda_i - i + j - 1}(x_1, \ldots, x_{k}|\tau^{1-j}y).
    \end{multline*}
    Notice that the factor 
    \[
    x_k - y_{k + \lambda_i - i} = x_k - y_{k + c(\square)}
    \]
    corresponds exactly to the deletion of the box on the border of $\lambda$ in the $i$-th row in the formula \eqref{ReverseBranching}. Thus, we can write, using the linearity of the determinant, the initial Jacobi-Trudi determinant as the sum of the determinants with boxes deleted from the border of $\lambda$ and only one box deleted in the same row
    \[
    \det(h_{\lambda_{i} - i + j}(x_1, \ldots, x_{k - 1}|\tau^{1 - j}y)) = \sum\limits_{\lambda/\mu=(1^r)}(-1)^r \prod\limits_{\square \in \lambda / \mu}(x_k - y_{k + c(\square)}) s_{\mu}(x_1, \ldots, x_k|y),
    \]
    which finishes the proof of the proposition.
\end{proof}

\end{document}
